\chapter{Introduction}
\label{sec:introduction}

\section{Context}
\label{subsec:context}

Amyotrophic lateral sclerosis (\gls{ALS}) is a fatal neurodegenerative disease characterised by progressive degeneration of upper and lower motor neurons, leading to muscle weakness, dysarthria, dysphagia, and respiratory failure \cite{Feldman2022}. Median survival from symptom onset is 2--4 years for most patients, but the distribution is wide: approximately 10\% of patients survive beyond ten years, whilst others deteriorate within one to two years of diagnosis \cite{Feldman2022}. This heterogeneity in disease course is not well explained by any single clinical variable measured at first assessment, which makes individual prognosis a persistent and clinically consequential open problem.

Functional status in ALS is standardly quantified using the \gls{ALSFRS-R}, a twelve-item rating scale covering speech, salivation, swallowing, handwriting, cutting food, dressing, turning in bed, walking, climbing stairs, and three respiratory domains \cite{cedarbaum1999alsfrs}. Each item is scored from 0 to 4, yielding a maximum total of 48. The rate of \gls{ALSFRS-R} decline over time --- the functional progression slope --- is the clinical marker most commonly used to classify patients into rapid and slow progressors, to stratify enrolment in clinical trials, and to guide decisions about care escalation. However, the correlation between a patient's baseline \gls{ALSFRS-R} score and their subsequent slope is only moderate, and predicting the slope from a single baseline assessment is not reliably achieved by current clinical scoring systems.

The Pooled Resource Open-Access ALS Clinical Trials (\gls{PRO-ACT}) database aggregates longitudinal records from 23 ALS clinical trials, covering approximately 10,000 subjects across multiple countries \cite{PROACT, kueffner2019}. After standard clinical exclusions and baseline alignment, working cohorts of approximately 1,000--2,000 subjects are obtainable, depending on the target being studied and the completeness requirements imposed on the feature set. \gls{PRO-ACT} provides demographic, functional rating, vital signs, pulmonary function, laboratory, and treatment data, and served as the primary dataset for an international ALS progression prediction challenge \cite{kueffner2019}. Both studies reported in this dissertation use \gls{PRO-ACT} as their data source.

\section{Motivation}
\label{subsec:motivation}

Accurate individual prognosis in ALS has concrete clinical consequences. For clinical trial enrolment, stratification by predicted progression rate reduces between-subject variance within treatment arms and can lower the sample sizes required to detect treatment effects. For palliative care planning, identifying patients whose \gls{ALSFRS-R} total is expected to decline sharply within three to six months can advance decisions on ventilatory support, enteral nutrition, and hospice referral. For survival prognosis, distinguishing patients likely to die within 24 months of symptom onset from those likely to survive longer determines the appropriate intensity of multidisciplinary intervention.

Building \gls{ML} models for ALS prognosis involves a set of technical problems that are not consistently addressed in the published literature. The first is temporal leakage. In longitudinal datasets, patient records from different time points of the same individual must be assigned to the same data partition; if a patient's early and later visits are split across training and test sets, the model learns from future observations about its own test subjects and inflates apparent performance \cite{Grollemund2019}. The second is class imbalance. In functional progression prediction, rapid progressors typically represent approximately 30\% of the cohort, a moderate imbalance for which the precision--recall curve is more informative than the \gls{ROC-AUC} curve. In survival prognosis, short survivors may account for 10--15\% of patients; at such ratios, optimising \gls{ROC-AUC} can yield classifiers that rank the minority class correctly in aggregate but detect almost none at the default 0.5 decision threshold. This behaviour --- high ranking performance coexisting with near-zero sensitivity --- is documented in this dissertation for a multi-layer perceptron classifier under severe imbalance and represents a material risk if such models were applied in clinical practice without threshold analysis. The third problem is interpretability: predictions without feature-level explanations cannot be evaluated against clinical domain knowledge and cannot support individual patient decisions.

A fourth issue is reproducibility. Most published ALS prognosis studies report cross-validation performance without a dedicated held-out test set \cite{Papaiz2022, Halbersberg2023}. Statistical comparisons between classifiers are informal or absent, and uncertainty quantification on test-set metrics is rarely provided. Papaiz et al.\ \cite{Papaiz2024} proposed a BalancedBagging ensemble strategy for survival prognosis in ALS using \gls{PRO-ACT} and demonstrated statistically validated improvements over single-classifier configurations for five of six tested classifiers, but their study did not include Bayesian hyperparameter optimisation, a held-out test set, or \gls{SHAP}-based interpretability. The second study in this dissertation replicates and extends that work, addressing all three missing components.

\section{Objectives}
\label{subsec:objectives}

This dissertation addresses ALS prognosis through two complementary studies, both using \gls{PRO-ACT} and both requiring leakage-free evaluation, class imbalance handling, and post-hoc interpretability.

\textbf{Study I: Functional Progression Prediction.} The objective is to build and evaluate a leakage-free classification pipeline for predicting rapid versus slow \gls{ALSFRS-R} functional progression at three-month and six-month horizons. The positive class is defined as the worst 30\% of \gls{ALSFRS-R} slopes, computed fold-wise on the training partition to prevent the threshold from encoding held-out information. Thirty-five baseline clinical features are derived from five \gls{PRO-ACT} data blocks: demographics and baseline \gls{ALSFRS-R}, vital signs, forced vital capacity (\gls{FVC}), treatment history, and laboratory values. Seven classifiers are tuned via Optuna Bayesian optimisation \cite{optuna2019} (60 trials per configuration, \gls{PR-AUC} objective) under 5-fold GroupKFold cross-validation with subject-level partition assignment to prevent patient leakage. An 80/20 subject-level holdout split is applied before tuning; the test set remains sealed until final evaluation. \gls{SHAP} \cite{lundberg2017shap} and \gls{LIME} \cite{ribeiro2016lime} are applied to the selected XGBoost model to attribute predictions at the feature level. A per-patient risk report generator produces outputs across three risk tiers. The primary evaluation metric is \gls{PR-AUC} at the six-month horizon, with the F1-optimal threshold used to report recall, precision, and F1 score.

\textbf{Study II: Survival Prognosis Replication and Extension.} The objective is to replicate the Papaiz et al.\ \cite{Papaiz2024} pipeline for predicting short survival (death within 24 months of symptom onset) and to extend it along four dimensions: a seventh classifier (LightGBM), ten individual class-imbalance resampling strategies yielding 70 model--strategy combinations, Bayesian hyperparameter optimisation with Optuna (50 trials per combination, \gls{ROC-AUC} objective), and a held-out test set withheld from all training decisions. The analysis is supplemented with bootstrap confidence intervals ($B = 2{,}000$ resamples) \cite{efron1993bootstrap} on test-set metrics, Wilcoxon signed-rank tests with Bonferroni correction for the BalancedBagging comparison \cite{wilcoxon1945}, decision threshold analysis using Youden's J statistic \cite{youden1950} and F1-optimal thresholds, and dual \gls{SHAP} attribution using TreeSHAP \cite{lundberg2020treeshap} for LightGBM and KernelSHAP for the multi-layer perceptron. A central finding is the decoupling between \gls{ROC-AUC} and sensitivity under severe imbalance (11.6\% minority class): the multi-layer perceptron achieves a test \gls{ROC-AUC} of 0.918 whilst detecting only 1 of 35 short survivors at the default threshold, whereas LightGBM with random oversampling achieves 0.923 \gls{ROC-AUC} and recalls 30 of 35 short survivors.

\section{Workflow}
\label{subsec:workflow}

Both studies follow a shared modelling philosophy: define a clinically grounded binary target, engineer features strictly from information available at the patient's baseline visit, apply a leakage-free train--test split before any modelling, tune hyperparameters using Bayesian search on the training partition only, and explain the final model using \gls{SHAP} feature attribution.

At the data level, both studies load \gls{PRO-ACT} from 16 source CSV files, apply clinical exclusions, and construct baseline feature vectors. Study I defines baseline as the first \gls{ALSFRS-R} assessment and extracts 35 clinical features from five data blocks. Study II follows the Papaiz et al.\ specification, defining baseline using the diagnosis visit and extracting 23 ordinal and binary features; it applies additional filtering criteria requiring at least two \gls{ALSFRS-R} assessments per patient, a known site of onset, and non-missing \gls{FVC} and height measurements.

At the modelling level, both studies apply a stratified 80/20 subject-level holdout split before any training. Class imbalance is handled differently by design: Study I uses classifier class-weight parameters in two modes (balanced and unbalanced), while Study II evaluates ten explicit resampling strategies ranging from random oversampling and \gls{SMOTE} \cite{chawla2002smote} to random undersampling and hybrid methods such as SMOTETomek. Both studies use Optuna with Tree-structured Parzen Estimator sampling for hyperparameter search. The cross-validation strategy also differs: Study I uses 5-fold GroupKFold with subject-level group assignment, preventing any patient's records from appearing in both training and validation folds; Study II uses RepeatedStratifiedKFold (5 folds $\times$ 3 repeats, 15 evaluations per configuration) following the Papaiz et al.\ protocol.

At the evaluation level, both studies report held-out test metrics with bootstrap confidence intervals. Study I uses \gls{PR-AUC} as the primary ranking metric, reflecting the 30\%/70\% class split for which the precision--recall curve is more informative than \gls{ROC-AUC}. Study II uses \gls{ROC-AUC} to align with the replication objective, but reports sensitivity as the primary threshold-dependent metric, because detecting short survivors is the clinically critical task and \gls{ROC-AUC} alone does not guarantee adequate recall under severe imbalance.

\section{Document Structure}
\label{subsec:structure}

The remainder of this dissertation is structured as follows.

Chapter~\ref{sec:stateoftheart} reviews the clinical and computational background: ALS heterogeneity and prognosis, functional decline and survival outcomes, the \gls{PRO-ACT} dataset, machine learning approaches to ALS prognosis, class imbalance handling, \gls{XAI} methods, and reproducibility in clinical machine learning.

Chapter~\ref{ch:study1_methods} describes the methods for Study I: \gls{PRO-ACT} dataset construction, feature engineering across five data blocks, the leakage-free validation protocol, the Optuna hyperparameter optimisation configuration, class imbalance handling, threshold selection, and the explainability framework.

Chapter~\ref{ch:study1_results} reports and discusses the results of Study I: the feature block ablation analysis, cross-horizon comparison between the three-month and six-month models, held-out test evaluation with bootstrap confidence intervals, calibration, confusion matrices, \gls{SHAP} and \gls{LIME} explanations, and the per-patient risk report.

Chapter~\ref{ch:study2_methods} describes the methods for Study II: the replication design and its relationship to Study I, the 70-combination experimental grid, the BalancedBagging comparison protocol, the Optuna optimisation configuration for the replication setting, and the dual \gls{SHAP} attribution approach.

Chapter~\ref{ch:study2_results} reports and discusses the results of Study II: cross-validation performance across 70 model--strategy combinations, held-out test results with bootstrap confidence intervals, the \gls{ROC-AUC}--sensitivity decoupling phenomenon, decision threshold analysis, the BalancedBagging statistical evaluation, and \gls{SHAP}-based feature interpretation for LightGBM and the multi-layer perceptron.

Chapter~\ref{ch:conclusion} synthesises findings from both studies, discusses cross-study methodological lessons, identifies limitations, and outlines directions for future work.